\section{What formation would have to add}
\label{sec:22.1}
A disposition is a \emph{candidate for formation} when the best available causal account of what the system did is a rationale that stayed active across contexts, rather than retrieval or compliance with whatever signal was present. The standard is comparative on purpose: a capable enough imitation reproduces any signature an evaluator can observe, so what an evaluation establishes is which explanation the conduct supports better.

Formation keeps the three conditions a refusal has to meet and adds others. Met in the evaluator's own vocabulary with nothing at stake, the three establish moral competence. Candidate formation is the same three held under adverse incentives, trainer pressure, a cost real inside the evaluation, and intervening updates, on cases the trainers didn't write.

None of the three carries the inference alone. Reaching new cases without holding under pressure may be broad imitation — a pattern extended that never had to be held against anything. Holding under pressure without lifting when the reason is gone may be rigidity — a refusal surviving every pressure including a defeating fact. Lifting without reaching new cases may be a trigger with an exception list, since lifting on a listed fact is as learnable as refusing on a listed form.
